\begin{figure}
  \centering
  \begin{tikzpicture}[x=1cm,y=1cm]
    \foreach \hamPanelShift/\hamStripWidth/\hamPanelTitle in {0/.35/{Narrow strip},4.25/1.65/{Corners removed},8.5/3.15/{Full rectangle}} {
    \begin{scope}[xshift=\hamPanelShift cm]
      \node[ham/note] at (0,2.15) {\hamPanelTitle};
      \begin{scope}
        \clip (-1.55,-1.05) rectangle (1.55,1.05);
        \fill[hamteal!18] (-4,-4-\hamStripWidth)--(4,4-\hamStripWidth)--(4,4+\hamStripWidth)--(-4,-4+\hamStripWidth)--cycle;
        \draw[hamteal,densely dashed] (-4,-4-\hamStripWidth)--(4,4-\hamStripWidth);
        \draw[hamteal,densely dashed] (-4,-4+\hamStripWidth)--(4,4+\hamStripWidth);
      \end{scope}
      \draw[hamblue] (-1.55,-1.05) rectangle (1.55,1.05);
      \draw[ham/arrow] (-1.8,0)--(1.8,0) node[right] {$x$};
      \draw[ham/arrow] (0,-1.3)--(0,1.4) node[above] {$y$};
      \node[font=\small,anchor=north] at (-1.55,-1.1) {$-a$};
      \node[font=\small,anchor=north] at (1.55,-1.1) {$a$};
      \node[font=\small,anchor=east] at (-1.6,1.05) {$b$};
      \node[font=\small,anchor=east] at (-1.6,-1.05) {$-b$};
    \end{scope}
    }
    \node[font=\small] at (0,-1.95) {$0<c<a-b$};
    \node[font=\small] at (4.25,-1.95) {$a-b\le c<a+b$};
    \node[font=\small] at (8.5,-1.95) {$c\ge a+b$};
  \end{tikzpicture}
  \caption[The rectangle and strip in the fourth-order calculation]{The area
    $F(a,b,c)$ in the spatial-white order-four calculation: the shaded part is
    $[-a,a]\times[-b,b]\cap\{|x-y|<c\}$, and the dashed lines are the edges
    $|x-y|=c$ of the strip. The drawing uses $a>b>0$; here $a,b,c$
    are time-increment parameters, not covariance exponents. The three geometric
    cases explain the positive-part formula, without replacing the separate
    pairing and time-order calculations. See
    \ProofRef{Proposition}{P:white-four}.}
  \label{Fig:white-strip}
\end{figure}
